
\chapter{2019年上海卷（秋）}

\section{填空题}

\begin{problem}
已知集合 \(A=(-\infty,3)\)，\(B=(2,+\infty)\)，则 \(A\cap B=\)\fillinblank{}.

\end{problem}

\begin{answer}
\((2,3)\)
\end{answer}

\begin{solution}
由交集定义可知，
\(A\cap B=(2,3)\).
故答案为 \((2,3)\).
\end{solution}

\begin{problem}
已知 \(z\in\mathbb{C}\) 且满足 \(\frac1z-5=\i\)，求 \(z=\)\fillinblank{}.

\end{problem}

\begin{answer}
\(\frac{5-\i}{26}\)
\end{answer}

\begin{solution}
由 \(\frac1z-5=\i\)，得 \(\frac1z=5+\i\)，所以 \(z=\frac1{5+\i}=\frac{5-\i}{(5+\i)(5-\i)}=\frac{5-\i}{26}\).
故答案为 \(\frac{5-\i}{26}\).
\end{solution}

\begin{problem}
已知向量 \(\bs{a}=(1,0,2)\)，\(\bs{b}=(2,1,0)\)，则 \(\bs{a}\) 与 \(\bs{b}\) 的夹角为\fillinblank{}.

\end{problem}

\begin{answer}
\(\arccos\frac25\)
\end{answer}

\begin{solution}
由向量夹角公式，\(\cos\theta=\frac{\bs{a}\cdot\bs{b}}{|\bs{a}||\bs{b}|}=\frac{1\cdot2+0\cdot1+2\cdot0}{\sqrt{1^2+0^2+2^2}\sqrt{2^2+1^2+0^2}}=\frac25\).
故夹角为 \(\arccos\frac25\).
\end{solution}

\begin{problem}
已知二项式 \((2x+1)^5\)，则展开式中含 \(x^2\) 项的系数为\fillinblank{}.

\end{problem}

\begin{answer}
\(40\)
\end{answer}

\begin{solution}
由二项式展开式通项公式，\(T_{r+1}=C_5^r(2x)^{5-r}\). 令 \(5-r=2\)，得 \(r=3\). 因此含 \(x^2\) 项的系数为 \(C_5^3\cdot2^2=40\).
故答案为 \(40\).
\end{solution}

\begin{problem}
已知 \(x\ge0\)，\(y\ge0\)，\(x+y\le2\)，求 \(z=2x-3y\) 的最小值为\fillinblank{}.

\begin{center}
\bitmapinclude[max width=6.0cm,max totalheight=4.2cm]{img/2019/shanghai/q5_fig1.png}
\end{center}

\end{problem}

\begin{answer}
\(-6\)
\end{answer}

\begin{solution}
由约束条件作出可行域如图所示，目标函数 \(z=2x-3y\) 在可行域上的最小值出现在点 \((0,2)\). 代入得 \(z_{\min}=2\times0-3\times2=-6\).
故答案为 \(-6\).
\end{solution}

\begin{problem}
已知函数 \(f(x)\) 周期为 \(1\)，且当 \(0<x\le1\) 时，\(f(x)=-\log_2x\)，则 \(f\!\left(\frac32\right)=\)\fillinblank{}.

\end{problem}

\begin{answer}
\(1\)
\end{answer}

\begin{solution}
因为 \(f(x)\) 的周期为 \(1\)，所以 \(f\!\left(\frac32\right)=f\!\left(\frac12\right)\). 又 \(0<\frac12\le1\)，故 \(f\!\left(\frac12\right)=-\log_2\frac12=1\).
故答案为 \(1\).
\end{solution}

\begin{problem}
若 \(x\)，\(y\in\R^+\)，且 \(\frac1x+2y=3\)，则 \(\frac yx\) 的最大值为\fillinblank{}.

\end{problem}

\begin{answer}
\(\frac98\)
\end{answer}

\begin{solution}
由 \(\frac1x+2y=3\)，设 \(t=\frac1x>0\)，则 \(t+2y=3\)，\(y=\frac{3-t}{2}\). 于是 \(\frac yx=ty=\frac{t(3-t)}2=-\frac12\left(t-\frac32\right)^2+\frac98\).
故 \(\frac yx\) 的最大值为 \(\frac98\).
\end{solution}

\begin{problem}
已知数列 \(\{a_n\}\) 前 \(n\) 项和为 \(S_n\)，且满足 \(S_n+a_n=2\)，则 \(S_5=\)\fillinblank{}.

\end{problem}

\begin{answer}
\(\frac{31}{16}\)
\end{answer}

\begin{solution}
由 \(S_n+a_n=2\)，\(S_{n-1}+a_{n-1}=2\) 以及 \(S_n=S_{n-1}+a_n\)，得 \(S_{n-1}+2a_n=2\)，\(S_{n-1}+a_{n-1}=2\). 相减可得 \(2a_n=a_{n-1}\).
因此 \(\{a_n\}\) 是公比为 \(\frac12\) 的等比数列. 又 \(S_1+a_1=2\) 且 \(S_1=a_1\)，所以 \(a_1=1\).
于是 \(S_5=1+\frac12+\frac14+\frac18+\frac1{16}=\frac{31}{16}\).
故答案为 \(\frac{31}{16}\).
\end{solution}

\begin{problem}
过 \(y^2=4x\) 的焦点 \(F\) 并垂直于 \(x\) 轴的直线分别与 \(y^2=4x\) 交于 \(A\)，\(B\)，\(A\) 在 \(B\) 上方，\(M\) 为抛物线上一点，\(\overrightarrow{OM}=\lambda\overrightarrow{OA}+(\lambda-2)\overrightarrow{OB}\)，则 \(\lambda=\)\fillinblank{}.

\end{problem}

\begin{answer}
\(3\)
\end{answer}

\begin{solution}
抛物线 \(y^2=4x\) 的焦点为 \(F(1,0)\). 过 \(F\) 且垂直于 \(x\) 轴的直线为 \(x=1\)，与抛物线交于 \(A(1,2)\)，\(B(1,-2)\). 设 \(M(x,y)\)，则 \(\overrightarrow{OM}=\lambda\overrightarrow{OA}+(\lambda-2)\overrightarrow{OB}=\lambda(1,2)+(\lambda-2)(1,-2)=(2\lambda-2,4)\)，所以 \(M(2\lambda-2,4)\). 又 \(M\) 在抛物线上，故 \(4^2=4(2\lambda-2)\)，
解得 \(\lambda=3\). 故答案为 \(3\).
\end{solution}

\begin{problem}
某三位数密码锁，每位数字在 \(0\sim9\) 数字中选取，其中恰有两位数字相同的概率是\fillinblank{}.

\end{problem}

\begin{answer}
\(\frac{27}{100}\)
\end{answer}

\begin{solution}
三位数字共有 \(10^3=1000\) 种取法.
恰有两位数字相同的情形数为 \(C_{10}^1\cdot C_3^2\cdot C_9^1=10\cdot3\cdot9=270\). 所以所求概率为 \(\frac{270}{1000}=\frac{27}{100}\).
故答案为 \(\frac{27}{100}\).
\end{solution}

\begin{problem}
已知数列 \(\{a_n\}\) 满足 \(a_n<a_{n+1}\)（\(n\in\N^*\)），\(P_n(n,a_n)\) 在双曲线 \(\frac{x^2}{6}-\frac{y^2}{2}=1\) 上，则 \(\lim\limits_{n\to\infty}|P_nP_{n+1}|=\)\fillinblank{}.

\end{problem}

\begin{answer}
\(\frac{2\sqrt3}{3}\)
\end{answer}

\begin{solution}
由 \(P_n(n,a_n)\) 在双曲线 \(\frac{x^2}{6}-\frac{y^2}{2}=1\) 上，得
 \(a_n=\sqrt{\frac{n^2}{3}-2}\). 因为 \(a_n<a_{n+1}\)，所以取上支.
于是 \(|P_nP_{n+1}|=\sqrt{1+(a_{n+1}-a_n)^2}\). 当 \(n\to\infty\) 时，\(a_{n+1}-a_n\to\frac1{\sqrt3}\)，故 \(\lim\limits_{n\to\infty}|P_nP_{n+1}|=\sqrt{1+\frac13}=\frac{2\sqrt3}{3}\).
故答案为 \(\frac{2\sqrt3}{3}\).
\end{solution}

\begin{problem}
已知 \(f(x)=\left|\frac{2}{x-1}-a\right|\)（\(x>1\)，\(a>0\)），若 \(a=a_0\)，\(f(x)\) 与 \(x\) 轴交点为 \(A\)，\(f(x)\) 为曲线 \(L\)，在 \(L\) 上任意一点 \(P\)，总存在一点 \(Q\)（\(P\) 异于 \(A\)）使得 \(AP\perp AQ\) 且 \(|AP|=|AQ|\)，则 \(a_0=\)\fillinblank{}.

\end{problem}

\begin{answer}
\(\sqrt2\)
\end{answer}

\begin{solution}
由 \(f(x)=0\) 可得曲线 \(L\) 与 \(x\) 轴交于 \(A\left(1+\frac{2}{a},0\right)\). 将坐标原点平移到 \(A\) 处，记点 \(P\) 的坐标为 \((u,v)\). 则曲线方程可写为 \(v=\frac{a|u|}{u+\frac2a}\).

先取 \(u>0\). 此时 \(v=\frac{au}{u+\frac2a}\). 因为 \(AP\perp AQ\) 且 \(|AP|=|AQ|\)，所以可取 \(Q(-v,u)\). 由 \(Q\in L\)，得 \(u=\frac{av}{\frac2a-v}\). 代入 \(v=\dfrac{au}{u+\frac2a}\)，得到 \(u=\frac{a\cdot \dfrac{au}{u+\frac2a}}{\frac2a-\dfrac{au}{u+\frac2a}}=\frac{a^2u}{\frac2a\left(u+\frac2a\right)-au}\). 因为 \(u>0\)，所以 \(\frac2a\left(u+\frac2a\right)-au=a^2\).
该式对任意 \(u>0\) 都成立，因此比较系数可得
 \(\frac2a-a=0\)，从而 \(a^2=2\).
又 \(a>0\)，所以
 \(a_0=\sqrt2\).
故答案为 \(\sqrt2\).
\end{solution}
\section{单选题}

\begin{problem}
已知直线方程 \(2x-y+c=0\) 的一个方向向量 \(d\) 可以是（\quad）

\choices
    {\((2,-1)\)}
    {\((2,1)\)}
    {\((-1,2)\)}
    {\((1,2)\)}

\end{problem}

\begin{answer}
D
\end{answer}

\begin{solution}
直线 \(2x-y+c=0\) 的法向量为 \((2,-1)\)，所以它的一个方向向量可以是 \((1,2)\). 故选 D.
\end{solution}

\begin{problem}
一个直角三角形的两条直角边长分别为 \(1\) 和 \(2\)，将该三角形分别绕其两个直角边旋转得到的两个圆锥的体积之比为（\quad）

\choices
    {\(1\)}
    {\(2\)}
    {\(4\)}
    {\(8\)}

\end{problem}

\begin{answer}
B
\end{answer}

\begin{solution}
绕长为 \(1\) 的直角边旋转时，圆锥的底面半径为 \(2\)，高为 \(1\)，体积为
\(V_1=\frac13\pi\cdot2^2\cdot1=\frac{4\pi}{3}\).
绕长为 \(2\) 的直角边旋转时，圆锥的底面半径为 \(1\)，高为 \(2\)，体积为
\(V_2=\frac13\pi\cdot1^2\cdot2=\frac{2\pi}{3}\).
所以体积之比为
\(\frac{V_1}{V_2}=2\).
故选 B.
\end{solution}

\begin{problem}
已知 \(\omega\in\R\)，函数 \(f(x)=(x-6)^2\sin(\omega x)\)，存在常数 \(a\in\R\)，使得 \(f(x+a)\) 为偶函数，则 \(\omega\) 可能的值为（\quad）

\choices
    {\(\frac{\pi}{2}\)}
    {\(\frac{\pi}{3}\)}
    {\(\frac{\pi}{4}\)}
    {\(\frac{\pi}{5}\)}

\end{problem}

\begin{answer}
C
\end{answer}

\begin{solution}
要使 \(f(x+a)\) 为偶函数，先看前面的平方项：
\(f(x+a)=(x+a-6)^2\sin[\omega(x+a)]\). 若 \(f(x+a)\) 为偶函数，则需 \(a=6\)，这样前一部分变为 \(x^2\)，是偶函数.
此时只需 \(\sin[\omega(x+6)]\) 也是偶函数. 由 \(\sin[\omega(-x+6)]=\sin[\omega(x+6)]\) 对一切 \(x\) 成立，可得 \(\cos6\omega=0\). 所以 \(6\omega=\frac{\pi}{2}+k\pi\)（\(k\in\Z\)），即 \(\omega=\frac{(2k+1)\pi}{12}\).
在选项中只有 \(\frac{\pi}{4}\) 符合，故选 C.
\end{solution}

\begin{problem}
已知 \(\tan\alpha\cdot\tan\beta=\tan(\alpha+\beta)\).

\circled{1} 存在 \(\alpha\) 在第一象限，角 \(\beta\) 在第三象限；

\circled{2} 存在 \(\alpha\) 在第二象限，角 \(\beta\) 在第四象限.

\choices
    {\(\circled{1}\circled{2}\)均正确}
    {\(\circled{1}\circled{2}\)均错误}
    {\(\circled{1}\)对，\(\circled{2}\)错}
    {\(\circled{1}\)错，\(\circled{2}\)对}

\end{problem}

\begin{answer}
D
\end{answer}

\begin{solution}
设
\(x=\tan\alpha\)，\(y=\tan\beta\). 原式化为 \(xy=\frac{x+y}{1-xy}\)，即 \(x^2y^2+(1-x)y+x=0\).

对于 \(\circled{1}\)，若 \(\alpha\) 在第一象限，则 \(x>0\). 把上式看作关于 \(y\) 的二次方程，
 \(\Delta=(1-x)^2-4x^3\). 当 \(x>0\) 时，\(\Delta\) 随 \(x\) 增大而减小；且若 \(x>1\)，则 \(\Delta<0\)，无实根；若 \(0<x<1\)，则两根之和 \(\frac{x-1}{x^2}<0\)，两根乘积 \(\frac1x>0\)，
故两根都为负数，不可能对应第三象限的 \(\beta\)（此时应有 \(y>0\)）. 所以 \(\circled{1}\) 错.

对于 \(\circled{2}\)，若 \(\alpha\) 在第二象限，则 \(x<0\). 这时
 \(\Delta=(1-x)^2-4x^3>0\)，且两根乘积 \(\frac1x<0\)，所以方程有一正一负两根，故可以取 \(y<0\)，即 \(\beta\) 在第四象限. 故 \(\circled{2}\) 对.

所以选 D.
\end{solution}
\section{解答题}

\begin{problem}
如图，在长方体 \(ABCD-A_1B_1C_1D_1\) 中，\(M\) 为 \(BB_1\) 上一点，已知 \(BM=2\)，\(AD=4\)，\(CD=3\)，\(AA_1=5\).

\begin{enumerate}
\item 求直线 \(A_1C\) 与平面 \(ABCD\) 的夹角；
\item 求点 \(A\) 到平面 \(A_1MC\) 的距离.
\end{enumerate}

\begin{center}
\bitmapinclude[max width=6.0cm,max totalheight=4.8cm]{img/2019/shanghai/q17_fig1.png}
\end{center}
\end{problem}

\begin{solution}
（1）依题意，\(AA_1\perp\) 面 \(ABCD\)，连接 \(AC\)，则 \(A_1C\) 与平面 \(ABCD\) 所成夹角为 \(\angle A_1CA\). 因为 \(AA_1=5\)，\(AC=\sqrt{3^2+4^2}=5\)，所以 \(\triangle A_1CA\) 为等腰直角三角形，\(\angle A_1CA=\frac{\pi}{4}\). 所以直线 \(A_1C\) 与平面 \(ABCD\) 的夹角为 \(\frac{\pi}{4}\).

（2）法一（空间向量）：如图建立坐标系，则
 \(A(0,0,0)\)，\(C(3,4,0)\)，\(A_1(0,0,5)\)，\(M(3,0,2)\). 于是 \(\overrightarrow{AC}=(3,4,0)\)，\(\overrightarrow{A_1C}=(3,4,-5)\)，\(\overrightarrow{MC}=(0,4,-2)\).
设平面 \(A_1MC\) 的法向量为 \(\bs{n}=(x,y,z)\)，则
\examdisplaycases{}{
\bs{n}\cdot\overrightarrow{A_1C}=3x+4y-5z=0,\\
\bs{n}\cdot\overrightarrow{MC}=4y-2z=0
}
解得可取 \(\bs{n}=(2,1,2)\). 因此点 \(A\) 到平面 \(A_1MC\) 的距离为 \(d=\frac{\left|\overrightarrow{AC}\cdot\bs{n}\right|}{\left|\bs{n}\right|}=\frac{3\times2+4\times1}{\sqrt{2^2+1^2+2^2}}=\frac{10}{3}\).

法二（等体积法）：利用 \(V_{A-A_1MC}=V_{C-A_1AM}\) 求解. 求 \(S_{\triangle A_1MC}\) 时，需要求出三边长（不是特殊三角形），利用 \(S_{\triangle}=\frac12 ab\sin C\) 求解.
\end{solution}

\begin{problem}
已知 \(f(x)=ax+\frac1{x+1}\)（\(a\in\R\)）.

\begin{enumerate}
\item 当 \(a=1\) 时，求不等式 \(f(x)+1<f(x+1)\) 的解集；
\item 若 \(x\in[1,2]\) 时，\(f(x)\) 有零点，求 \(a\) 的范围.
\end{enumerate}

\end{problem}

\begin{solution}
（1）当 \(a=1\) 时，\(f(x)=x+\frac1{x+1}\).
代入原不等式：
 \(x+\frac1{x+1}+1<x+1+\frac1{x+2}\)，
即
 \(\frac1{x+1}<\frac1{x+2}\).
移项通分：\(\frac1{(x+1)(x+2)}<0\)，
解得 \(-2<x<-1\).

（2）依题意，\(f(x)=ax+\frac1{x+1}=0\) 在 \(x\in[1,2]\) 上有解.
参数分离得
 \(a=-\frac1{x(x+1)}\)，
即求 \(g(x)=-\frac1{x(x+1)}\) 在 \(x\in[1,2]\) 上的值域.
因为 \(x(x+1)\) 在 \(x\in[1,2]\) 上单调递增，且 \(x(x+1)\in[2,6]\)，所以
 \(-\frac1{x(x+1)}\in\left[-\frac16,-\frac12\right]\).
故 \(a\in\left[-\frac12,-\frac16\right]\).
\end{solution}

\begin{problem}
如图，\(A-B-C\) 为海岸线，\(AB\) 为线段，\(BC\) 为四分之一圆弧，\(BD=39.2\text{ km}\)，\(\angle BDC=22^\circ\)，\(\angle CBD=68^\circ\)，\(\angle BDA=58^\circ\).

\begin{enumerate}
\item 求弧 \(BC\) 长度；
\item 若 \(AB=40\text{ km}\)，求 \(D\) 到海岸线 \(A-B-C\) 的最短距离（精确到 \(0.001\text{ km}\)）.
\end{enumerate}

\begin{center}
\bitmapinclude[max width=6.0cm,max totalheight=4.8cm]{img/2019/shanghai/q19_fig1.png}
\end{center}
\end{problem}

\begin{solution}
（1）依题意，\(|BC|=BD\cdot\sin22^\circ\). 弧 \(BC\) 所在圆的半径 \(R=|BC|\cdot\sin\frac{\pi}{4}\). 弧 \(BC\) 长度为 \(\frac{\pi}{2}R=\frac{\pi}{2}\cdot |BC|\cdot\frac{\sqrt2}{2}=\frac{\sqrt2}{4}\times 3.141\times 39.2\times \sin22^\circ=16.310\text{ km}\).

（2）根据正弦定理，\(\frac{BD}{\sin A}=\frac{AB}{\sin58^\circ}\)，得 \(\sin A=\frac{39.2}{40}\sin58^\circ=0.831\)，\(A=56.2^\circ\). \(\angle ABD=180^\circ-56.2^\circ-58^\circ=65.8^\circ\). 所以 \(DH=BD\cdot\sin\angle ABD=35.752\text{ km}<CD=36.346\text{ km}\).
因此 \(D\) 到海岸线最短距离为 \(35.752\text{ km}\).
\end{solution}

\begin{problem}
已知椭圆 \(\frac{x^2}{8}+\frac{y^2}{4}=1\)，\(F_1\)，\(F_2\) 为左，右焦点，直线 \(l\) 过 \(F_2\) 交椭圆于 \(A\)，\(B\) 两点.

\begin{enumerate}
\item 若 \(AB\) 垂直于 \(x\) 轴时，求 \(|AB|\)；
\item 当 \(\angle F_1AB=90^\circ\) 时，\(A\) 在 \(x\) 轴上方时，求 \(A\)，\(B\) 的坐标；
\item 若直线 \(AF_1\) 交 \(y\) 轴于 \(M\)，直线 \(BF_1\) 交 \(y\) 轴于 \(N\)，是否存在直线 \(l\)，使 \(S_{\triangle F_1AB}=S_{\triangle F_1MN}\)，若存在，求出直线 \(l\) 的方程；若不存在，请说明理由.
\end{enumerate}

\end{problem}

\begin{solution}
（1）依题意，\(F_2(2,0)\). 当 \(AB\perp x\) 轴时，可设 \(A(2,\sqrt2)\)，\(B(2,-\sqrt2)\)，所以 \(|AB|=2\sqrt2\).

（2）法一：由 \(\angle F_1AB=90^\circ\) 可得 \(\angle F_1AF_2=\frac{\pi}{2}\).
由焦点三角形面积公式，\(S_{\triangle F_1AF_2}=b^2\cdot\tan\frac{\theta}{2}=4\tan\frac{\pi}{4}=4\). 又 \(|F_1F_2|=2c=4\)，故 \(S_{\triangle F_1AF_2}=\frac12\cdot 2c\cdot y_A=2y_A=4\)，即 \(y_A=2\). 所以 \(A\) 在短轴端点，即 \(A(0,2)\). 直线 \(l_{AF}\)（即 \(l_{AB}\)）方程为 \(y=-x+2\).
联立
\examdisplaycases{}{
y=-x+2,\\
\frac{x^2}{8}+\frac{y^2}{4}=1
}
得 \(B\left(\frac83,-\frac23\right)\).

法二：设 \(A(x_1,y_1)\)，则 \(\angle F_1AF_2=\frac{\pi}{2}\)，所以 \(\overrightarrow{AF_1}\cdot\overrightarrow{AF_2}=(x_1+2,y_1)\cdot(x_1-2,y_1)=x_1^2-4+y_1^2=0\). 又 \(A\) 在椭圆上，满足 \(\frac{x_1^2}{8}+\frac{y_1^2}{4}=1\)，即 \(y_1^2=4\left( 1-\frac{x_1^2}{8} \right)\). 代入得 \(x_1^2-4+4\left( 1-\frac{x_1^2}{8} \right)=0\)，解得 \(x_1=0\)，故 \(A(0,2)\). \(B\) 点坐标求解同法一，得 \(B\left(\frac83,-\frac23\right)\).

（3）设直线 \(l\) 的方程为 \(x=my+2\)（\(k\) 不存在时不满足题意）.
设 \(A(x_1,y_1)\)，\(B(x_2,y_2)\)，\(M(0,y_3)\)，\(N(0,y_4)\)，联立
\examdisplaycases{}{
x=my+2,\\
\frac{x^2}{8}+\frac{y^2}{4}=1
}
得 \((m^2+2)y^2+4my-4=0\). 由韦达定理，\(y_1+y_2=-\frac{4m}{m^2+2}\)，\(y_1y_2=-\frac4{m^2+2}\). 又 \(S_{\triangle F_1AB}=\frac12|F_1F_2|\cdot|y_1-y_2|=2|y_1-y_2|\)，\(S_{\triangle F_1MN}=\frac12|F_1O|\cdot|y_3-y_4|=|y_3-y_4|\).
由直线 \(AF_1\) 方程 \(y=\frac{y_1}{x_1+2}(x+2)\)，得 \(M\) 纵坐标 \(y_3=\frac{2y_1}{x_1+2}\). 由直线 \(BF_1\) 方程 \(y=\frac{y_2}{x_2+2}(x+2)\)，得 \(N\) 纵坐标 \(y_4=\frac{2y_2}{x_2+2}\).
若 \(S_{\triangle F_1AB}=S_{\triangle F_1MN}\)，即 \(2|y_1-y_2|=|y_3-y_4|\). 而 \(\left| y_3-y_4 \right|=\left| \frac{2y_1}{x_1+2}-\frac{2y_2}{x_2+2} \right|=\left| \frac{2y_1}{my_1+4}-\frac{2y_2}{my_2+4} \right|=\left| \frac{8(y_1-y_2)}{(my_1+4)(my_2+4)} \right|\). 所以 \(|(my_1+4)(my_2+4)|=4\). 即 \(\left| m^2y_1y_2+4m(y_1+y_2)+16 \right|=4\). 代入韦达定理，得 \(\left| -\frac{4m^2}{m^2+2}+4m\cdot\left( -\frac{4m}{m^2+2} \right)+16 \right|=4\)，解得 \(m=\pm\sqrt3\). 因此存在直线 \(x+\sqrt3\,y-2=0\) 或 \(x-\sqrt3\,y-2=0\) 满足题意.
\end{solution}

\begin{problem}
数列 \(\{a_n\}\) 有 \(100\) 项，\(a_1=a\)，对任意 \(n\in[2,100]\)，存在 \(a_n=a_i+d\)，\(i\in[1,n-1]\)，若 \(a_k\) 与前 \(n\) 项中某一项相等，则称 \(a_k\) 具有性质 \(P\).

\begin{enumerate}
\item 若 \(a_1=1\)，求 \(a_4\) 可能的值；
\item 若 \(\{a_n\}\) 不为等差数列，求证：\(\{a_n\}\) 中存在满足性质 \(P\)；
\item 若 \(\{a_n\}\) 中恰有三项具有性质 \(P\)，这三项和为 \(C\)，使用 \(a\)，\(d\)，\(c\) 表示 \(a_1+a_2+\cdots+a_{100}\).
\end{enumerate}

\end{problem}

\begin{solution}
（1）\(a_4\) 可能的值为 \(3\)，\(5\)，\(7\).

（2）要证明 \(\{a_n\}\) 中存在满足性质 \(P\) 的项，只需证明其逆否命题：若数列 \(\{a_n\}\) 中不存在满足性质 \(P\) 的项，则 \(\{a_n\}\) 为等差数列.
由题设可知
\(a_1=a\).
当 \(n=2\) 时，存在 \(i\in[1,1]\)，使得
\(a_2=a_i+d=a_1+d=a+d\).
当 \(n=3\) 时，存在 \(i\in[1,2]\)，使得
\(a_3=a_i+d\).
当 \(i=1\) 时，
\(a_3=a_1+d=a_2\)，
满足性质 \(P\)，不成立.
当 \(i=2\) 时，
\(a_3=a_2+d=a+2d\)，
此时
\(a_4=a_i+d\)（\(i=1\)，\(2\)，\(3\)）.
当 \(i=1\) 时，
\(a_4=a_1+d=a+d=a_2\)，
满足性质 \(P\)，不成立.

当 \(i=2\) 时，
\(a_4=a_2+d=(a+d)+d=a+2d=a_3\)，
满足性质 \(P\)，不成立.

所以
\(a_4=a_3+d=a+3d\).
以此类推，\(a_n=a_i+d\)（\(i\in[1,n-1]\)），其中 \(a_n=a_i+d\)（\(i\in[1,n-2]\)）时不成立，只有 \(i=n-1\)，即 \(a_n=a_{n-1}+d\) 成立，所以 \(\{a_n\}\) 为等差数列.
因此可得：\(\{a_n\}\) 不为等差数列时，\(\{a_n\}\) 中存在满足性质 \(P\) 的项.

（3）将数列中具有性质 \(P\) 的三项去掉，形成一个新数列 \(\{b_n\}\).
\(b_1=a_1=a\).
当 \(n\in[2,97]\) 时，
\(b_n=a_i+d\)，\(i\in[1,n-1]\)，
且 \(\{b_n\}\) 中不存在满足性质 \(P\) 的项. 根据（2），\(\{b_n\}\) 为等差数列，所以
\(b_1+b_2+\cdots+b_{97}=97b_1+\frac{97\times96}{2}d=97a+4656d\).
又因为三项去掉和为 \(c\)，所以
\(a_1+a_2+\cdots+a_{100}=97a+4656d+c\).
\end{solution}
